\documentclass[10pt,conference]{IEEEtran}
\IEEEoverridecommandlockouts

\usepackage{amsmath,amssymb,amsfonts}
\usepackage{algorithmic}
\usepackage{booktabs}
\usepackage{cite}
\usepackage{graphicx}
\usepackage{microtype}
\usepackage{multirow}
\usepackage{xcolor}
\usepackage{url}

\newcommand{\smoke}{\textsc{Smoke}}
\newcommand{\todoresult}[1]{\textcolor{red}{[FINAL CAMPAIGN REQUIRED: #1]}}

\begin{document}

\title{Deep Reinforcement Learning Control of UAV-Enabled ISAC Systems for Livestock Monitoring}

\author{
\IEEEauthorblockN{Oluwaseun Odusole\IEEEauthorrefmark{2},
Awanthi Jayasundara\IEEEauthorrefmark{2},
Engin Zeydan\IEEEauthorrefmark{1},
Madhusanka Liyanage\IEEEauthorrefmark{3}, and
Pasika Ranaweera\IEEEauthorrefmark{2}}
\IEEEauthorblockA{\IEEEauthorrefmark{1}CTTC/CERCA, Castelldefels, Spain}
\IEEEauthorblockA{\IEEEauthorrefmark{2}School of Electrical and Electronic Engineering,
University College Dublin, Ireland}
\IEEEauthorblockA{\IEEEauthorrefmark{3}NetsLab, School of Computer Science,
University College Dublin, Ireland}
}

\maketitle

\begin{abstract}
Large-area livestock monitoring requires sensing coverage, reliable data
delivery, and energy-aware aerial operation under time-varying conditions.
This paper develops a sequential deep reinforcement learning (DRL) controller
for a four-UAV integrated sensing and communication (ISAC) system simulated in
ns-3 with 5G New Radio connectivity. Unlike configuration-level surrogate
search, the proposed controller observes sensing, communication, energy,
mobility, herd, and mission-time information every 10~s and changes sensing
intervals, target speeds, or target altitudes during the same episode. A
synchronous Gymnasium--ns-3 interface enforces one state--action--transition
handshake per control window. The policy is learned with Double Deep
Q-Networks (Double DQN), experience replay, a separate target network, and
epsilon-greedy exploration. The implementation exposes a 50-dimensional state
and 25 discrete actions and computes throughput, delay, loss, detection,
localisation, and energy over non-overlapping control windows. Unit tests
verify the Bellman target and replay path, while an ns-3 integration test
confirms that sensing, speed, and altitude actions modify live simulation
state. A short three-episode experiment is reported only as an implementation
smoke test; statistically powered performance claims are reserved for the
predeclared multi-seed evaluation protocol provided in this paper.
\end{abstract}

\begin{IEEEkeywords}
deep reinforcement learning, Double DQN, integrated sensing and communication,
UAV, livestock monitoring, 5G NR, ns-3, Gymnasium, O-RAN
\end{IEEEkeywords}

\section{Introduction}

Continuous livestock observation over a large farm is difficult when fixed
sensors have incomplete coverage and rural connectivity is limited. Unmanned
aerial vehicles (UAVs) can move their sensors towards poorly observed areas
and relay measurements through a cellular network. Integrated sensing and
communication (ISAC) further allows a common aerial platform to sense targets
and deliver sensing reports. These benefits introduce a sequential control
problem: a short sensing interval may improve observation frequency but consume
additional processing and radio energy; increased speed can improve spatial
coverage but increases propulsion energy; and altitude changes affect both
geometry and radar range.

The earlier version of this study used a Random Forest (RF) regressor to map a
complete, fixed simulator configuration to its final objective value. That is
a useful surrogate-assisted random-search benchmark, but it is not
reinforcement learning: it contains neither state transitions nor temporally
ordered actions, Bellman updates, experience replay, or a learned policy that
acts during an episode. The present work corrects that methodological issue by
formulating and implementing an online sequential Markov decision process
(MDP). Deep Q-Networks combine Q-learning with neural function approximation
\cite{mnih2015}, while Double DQN reduces the maximisation bias that can occur
when the same estimator both selects and evaluates the next action
\cite{hasselt2016}.

The principal contributions are:
\begin{itemize}
  \item a reproducible ns-3/Gymnasium environment in which every action changes
  the running UAV simulation and produces exactly one subsequent KPI window;
  \item a 50-dimensional state, 25-action control space, energy-aware reward,
  and separate termination and time-truncation conditions;
  \item a tested Double-DQN implementation with replay memory, online and
  target networks, Huber loss, and gradient clipping;
  \item window-specific communication KPIs derived from FlowMonitor deltas,
  rather than cumulative statistics incorrectly reused at every decision; and
  \item an evaluation protocol that separates training, validation, and final
  test seeds and compares frozen policies on identical test seeds.
\end{itemize}

\section{System Model}

\subsection{Scenario and UAV Roles}

The simulator contains 80 mobile ground targets in an $800\times800$~m farm,
four UAV user equipments, one 5G NR gNB, and a remote monitoring host. The UAVs
have surveillance, patrol, rapid-response, and strategic roles with distinct
speed and altitude limits. Cow positions are updated every 1~s. Each UAV
periodically performs monostatic radar-like sensing and transmits detections
and telemetry through the NR network.

UAV motion follows a correlated Gauss--Markov update,
\begin{align}
\mathbf v_i[n] &= \alpha_i\mathbf v_i[n-1]
 +(1-\alpha_i)\bar{\mathbf v}_i
 +\sqrt{1-\alpha_i^2}\,\mathbf w_i[n],\\
\mathbf p_i[n] &= \mathbf p_i[n-1]+\mathbf v_i[n]\Delta t_i,
\end{align}
where $\alpha_i$ is the memory coefficient and $\mathbf w_i[n]$ is a Gaussian
innovation. The controller changes a live target speed or altitude; the
mobility model then approaches the target while retaining correlated motion.

\subsection{Sensing and Cooperative Localisation}

For UAV $i$ and animal $k$, the three-dimensional slant range is
\begin{equation}
R_{ik}(t)=\|\mathbf p_i(t)-\mathbf p_k(t)\|_2.
\end{equation}
The ground target has $z=0$, so altitude is included in $R_{ik}$. This detail is
essential because the monostatic received power varies as $R^{-4}$:
\begin{equation}
P_{r,ik}(t)=\frac{P_tG^2\lambda^2\sigma}
{(4\pi)^3R_{ik}^4(t)}.
\end{equation}
A detection occurs when $10\log_{10}(P_{r,ik}/P_n)$ exceeds the configured
threshold, with stochastic misses close to the threshold. SNR-weighted fusion
combines simultaneous estimates:
\begin{equation}
\hat{\mathbf p}_k=
\frac{\sum_{i\in\mathcal U_k}w_{ik}\hat{\mathbf p}_{ik}}
{\sum_{i\in\mathcal U_k}w_{ik}},\qquad
w_{ik}=10^{\mathrm{SNR}_{ik}/10}.
\end{equation}
The sensing KPIs are detection probability $P_{\rm det}$ and localisation root
mean-square error $R_{\rm mse}$.

\subsection{Window Communication KPIs}

Let $\Delta B^{\rm rx}_{i,t}$, $\Delta N^{\rm rx}_{i,t}$, and
$\Delta N^{\rm tx}_{i,t}$ denote FlowMonitor counter differences for ISAC flow
$i$ during decision window $t$ of duration $\Delta_c$. The window metrics are
\begin{align}
T_t &= \frac{1}{|\mathcal F|}\sum_{i\in\mathcal F}
\frac{8\Delta B^{\rm rx}_{i,t}}{10^6\Delta_c},\\
D_t &= \frac{1}{|\mathcal F|}\sum_{i\in\mathcal F}
\frac{\Delta d_{i,t}}{\max(1,\Delta N^{\rm rx}_{i,t})},\\
L_t &= \frac{100}{|\mathcal F|}\sum_{i\in\mathcal F}
\frac{\Delta N^{\rm lost}_{i,t}}{\max(1,\Delta N^{\rm tx}_{i,t})}.
\end{align}
Only ISAC-report flows are included. Counter snapshots are stored after each
decision, preventing cumulative history from being counted as the current
transition.

\subsection{Energy Model}

The cumulative energy of UAV $i$ is
\begin{equation}
E_i=E_i^{\rm prop}+E_i^{\rm RF}+E_i^{\rm proc}.
\end{equation}
Propulsion power contains hover and velocity-dependent terms,
\begin{equation}
P_i^{\rm prop}=P_i^{\rm hover}+c_v\|\mathbf v_i\|^2,
\end{equation}
RF energy integrates transmit power, and per-sensing-step processing energy is
\begin{equation}
e_i^{\rm proc}=\kappa(c_0+c_1N_i^{\rm det})f_i^2.
\end{equation}
The controller receives the remaining battery fraction for every UAV and an
episode terminates if any cumulative energy reaches its budget.

\section{Sequential DRL Formulation}

\subsection{Markov Decision Process}

At $t=\Delta_c,2\Delta_c,\ldots$, the environment returns
$(\mathbf s_t,r_t,\mathtt{terminated},\mathtt{truncated})$. Table~\ref{tab:state}
defines the state. All values are scaled to bounded numerical ranges; signed
horizontal positions and velocities lie in $[-1,1]$, while the remaining
features lie in $[0,1]$.

\begin{table}[t]
\caption{DRL State Vector ($|\mathcal S|=50$)}
\label{tab:state}
\centering
\scriptsize
\begin{tabular}{lcc}
\toprule
State group & Features & Count\\
\midrule
Window KPIs & $P_{\rm det}$, RMSE, throughput, delay, loss & 5\\
Battery & Remaining fraction for four UAVs & 4\\
UAV kinematics & $(x,y,z,v_x,v_y,v_z)$ per UAV & 24\\
Herd summary & Centroid $(x,y)$ and spread $(\sigma_x,\sigma_y)$ & 4\\
Live controls & Four sensing intervals & 4\\
 & Four target speeds & 4\\
 & Four target altitudes & 4\\
Mission phase & Normalised simulation time & 1\\
\midrule
Total & & 50\\
\bottomrule
\end{tabular}
\end{table}

The discrete action set contains no-op plus six actions for each UAV role:
decrease/increase sensing interval by 0.1~s, decrease/increase target speed by
1~m/s, and decrease/increase target altitude by 5~m. Thus,
$|\mathcal A|=1+4\times6=25$. Each value is clipped to its physical role bound.

Using the same normalisation constants as the simulator objective, the window
utility is
\begin{equation}
J_t=3\tilde P_{{\rm det},t}+3\tilde T_t-2\widetilde{R}_{{\rm mse},t}
-\tilde D_t-\tilde L_t.
\end{equation}
Let $\Delta E_t$ be fleet energy consumed during the window and
$E_t^{\rm budget}$ the proportional window budget. The reward is
\begin{equation}
r_t=J_t-\lambda_E\frac{\Delta E_t}{\max(E_t^{\rm budget},\epsilon)}.
\label{eq:reward}
\end{equation}
Battery exhaustion sets \texttt{terminated}; reaching mission time sets
\texttt{truncated}. The mission phase in the state makes the finite horizon
observable.

\subsection{Synchronous ns-3--Gymnasium Handshake}

The C++ simulator opens a local TCP server. After a control window it sends one
newline-delimited JSON transition containing episode ID, step ID, state,
reward, flags, and raw KPIs. Python validates the IDs and 50-value observation,
sends one integer action with matching IDs, and blocks until the next window.
This lock-step design prevents stale actions, missing decisions, and ambiguous
credit assignment. Episode directories store the full command, KPI log,
summary, and simulator output and are immutable to prevent accidental
appending or overwriting. The Gymnasium API follows the standard
\texttt{reset}/\texttt{step} interface \cite{towers2024}.

\subsection{Double Deep Q-Network}

The online network $Q_\theta(\mathbf s,a)$ is a multilayer perceptron with two
256-unit ReLU hidden layers and 25 outputs. A replay memory stores
$(\mathbf s_t,a_t,r_t,\mathbf s_{t+1},d_t)$. For a sampled transition, Double
DQN chooses the next action with the online network and evaluates it with the
target network $Q_{\bar\theta}$:
\begin{align}
a^*&=\arg\max_a Q_\theta(\mathbf s_{t+1},a),\\
y_t&=r_t+\gamma(1-d_t)Q_{\bar\theta}(\mathbf s_{t+1},a^*).
\label{eq:ddqn}
\end{align}
The Huber loss between $Q_\theta(\mathbf s_t,a_t)$ and $y_t$ is minimised using
AdamW. Gradients are clipped to norm 10. The target weights are copied from the
online network every 1000 environment steps. Exploration decreases linearly
from $\epsilon=1.0$ to 0.05 over 30000 steps.

\begin{algorithmic}[1]
\STATE Initialise online network $Q_\theta$, target network
$Q_{\bar\theta}\leftarrow Q_\theta$, and replay memory $\mathcal D$
\FOR{each training episode}
  \STATE Reset ns-3 with a predeclared training seed and receive $\mathbf s_0$
  \WHILE{episode is active}
    \STATE Select random action with probability $\epsilon$; otherwise
    $a_t=\arg\max_a Q_\theta(\mathbf s_t,a)$
    \STATE Apply $a_t$ in live ns-3 and receive $(\mathbf s_{t+1},r_t,d_t)$
    \STATE Store transition in $\mathcal D$
    \STATE Sample a mini-batch and update $\theta$ using~\eqref{eq:ddqn}
    \STATE Periodically copy $\bar\theta\leftarrow\theta$
  \ENDWHILE
\ENDFOR
\end{algorithmic}

\section{Experimental Methodology}

\subsection{Simulation and Learning Parameters}

\begin{table}[t]
\caption{Implemented Default Parameters}
\label{tab:params}
\centering
\scriptsize
\begin{tabular}{ll}
\toprule
Parameter & Value\\
\midrule
UAVs; cows; field & 4; 80; $800\times800$~m\\
Mission; control interval & 300~s; 10~s\\
Carrier; NR bandwidth & 3.5~GHz; 20~MHz\\
UE/gNB transmit power & 23/38~dBm\\
Sensing transmit power/gain & 30~dBm/20~dB\\
SNR threshold; position noise & 30~dB; 5~m\\
Initial UAV altitudes & 60, 80, 100, 120~m\\
State/action dimensions & 50/25\\
Hidden layers & 256, 256 ReLU\\
$\gamma$; learning rate & 0.99; $3\times10^{-4}$\\
Replay/batch size & 100000/128\\
Learning starts; target update & 2000/1000 steps\\
Exploration & 1.0 to 0.05 over 30000 steps\\
\bottomrule
\end{tabular}
\end{table}

Table~\ref{tab:params} reports code defaults, correcting inconsistencies in the
earlier manuscript, which listed a 600~s duration and 15~dB threshold while the
available implementation used 300~s and 30~dB. Any non-default final setting
must be recorded by the generated per-episode command manifest.

\subsection{Baselines and Fair Comparison}

The comparison contains: (i) static/no-op control; (ii) uniformly random
actions with a recorded RNG seed; (iii) a transparent threshold rule; (iv) the
released 28-feature RF surrogate configuration, held static within an episode;
and (v) the frozen Double-DQN policy. The RF model is a benchmark and is not
described as RL. Its released feature set excludes UE and gNB transmit power,
although the earlier thesis counted them when referring to 30 variables.

Training, validation, and test seeds are disjoint. Hyperparameters and
checkpoint selection use validation seeds only. After selection, the network
is frozen and every method is evaluated on identical unseen test seeds. The
primary analysis reports paired per-seed differences, mean, standard
deviation, and 95\% confidence intervals. Multiple independent training seeds
should additionally quantify optimisation variability.

\section{Verification and Preliminary Results}

\subsection{Software Verification}

\begin{table}[t]
\caption{Completed Verification Tests}
\label{tab:tests}
\centering
\scriptsize
\begin{tabular}{p{0.49\linewidth}p{0.39\linewidth}}
\toprule
Test & Verified property\\
\midrule
Analytic Double-DQN target & Online argmax and target-network value; terminal masking\\
Replay-memory test & State/action/reward batch shapes and sampling path\\
Mock-environment learning & Greedy policy learns known preferred action 7\\
Protocol socket-pair tests & Episode/step IDs, action transfer, 50-state validation\\
ns-3 integration test & Live sensing, target speed, target altitude, and actual altitude change\\
Build and syntax checks & ns-3 executable, Python compilation, whitespace checks\\
\bottomrule
\end{tabular}
\end{table}

The ns-3 integration test completed in 66.9~s. It decreased the surveillance
sensing interval, increased its speed target, and increased its altitude target
across consecutive decision windows, then verified the corresponding state
changes. An inherited range defect was also corrected: the earlier code used
the UAV altitude at both range endpoints and therefore removed vertical
distance from radar sensing. The corrected calculation measures UAV-to-ground
slant range.

\subsection{End-to-End Smoke Test}

The post-fix smoke campaign used three 30~s training episodes, producing only
six transitions. Candidate checkpoints were evaluated on validation seed 7002;
all tied, and the first was frozen. Table~\ref{tab:smoke} reports one subsequent
matched test seed (9002). These data verify execution and file generation; they
are far too small for algorithm ranking.

\begin{table}[t]
\caption{Post-Fix Pipeline Smoke Test on One Unseen Seed}
\label{tab:smoke}
\centering
\scriptsize
\begin{tabular}{lrrrr}
\toprule
Policy & Return & $P_{\rm det}$ & RMSE (m) & Energy (J)$^a$\\
\midrule
Static & 2.6033 & 0.5609 & 9.4829 & 11239.26\\
Random & 2.7903 & 0.5808 & 8.9565 & 11239.32\\
Rule & 2.7215 & 0.5665 & 9.1766 & 11239.30\\
Double DQN & 2.5076 & 0.5655 & 9.4867 & 11239.20\\
\bottomrule
\multicolumn{5}{l}{\scriptsize $^a$Energy in the final 10~s window. One seed; no CI can be estimated.}
\end{tabular}
\end{table}

The smoke DQN does not outperform the baselines, nor should it be expected to:
its epsilon remained approximately 1.0 and it received six transitions, below
the default 2000-step learning-start threshold used for a real campaign. The
smoke run deliberately used a reduced threshold only to exercise gradient
updates. Consequently, Table~\ref{tab:smoke} must not appear as evidence that
DRL improves the system.

\subsection{Required Final Results}

Before submission, replace this subsection with automatically generated tables
from the full campaign. At minimum report:
\begin{itemize}
  \item learning curves across independent agent seeds with validation return;
  \item final static, RF, random, rule, and frozen-DQN results on the same
  unseen seeds;
  \item paired differences and 95\% confidence intervals for return,
  $P_{\rm det}$, RMSE, throughput, delay, loss, and energy;
  \item action frequencies and bound-hit rates to show what the policy learned;
  \item ablations removing energy, communication KPIs, and mobility actions;
  and
  \item wall-clock cost and sensitivity to control interval.
\end{itemize}

\todoresult{Insert the multi-seed final table and learning curves generated by
the documented full-campaign commands. Do not reuse the legacy 40/14-run RF
numbers as DRL results.}

\section{Reproducibility and Limitations}

The audited public repository contains one commit, no tags, and no original
CSV, model, or command manifest. The supplied thesis PDF contains no embedded
data. Therefore, the exact 300-plus-run thesis dataset cannot be reconstructed
from the available artifacts. New RF datasets are generated with a fixed
sampling seed and a saved configuration and command for each run. Every DRL
episode also records its command, seed, KPI windows, simulator log, and final
summary. Existing episode directories are never silently overwritten.

The simulator is a controlled abstraction rather than a field validation.
Radar returns, animal motion, energy, channel effects, and processing load are
modelled; hardware faults, weather, regulatory restrictions, and real animal
behaviour are not. The 50-feature state uses herd centroid and spread rather
than all animal positions, so the process may be partially observed. Future
work should compare recurrent agents, validate models against field data, and
deploy only a tested frozen policy through an O-RAN xApp with safety bounds and
a fallback controller.

\section{Conclusion}

This work replaces a mislabeled RF surrogate search with a genuine sequential
DRL controller for UAV-enabled ISAC livestock monitoring. The implementation
connects ns-3 and Gymnasium through a strict state--action handshake, applies
sensing and mobility actions to live objects, computes communication KPIs per
decision window, and trains a tested Double-DQN agent. Unit and integration
tests establish software correctness, and a small smoke campaign establishes
end-to-end operability. It does not establish a performance advantage. The
provided seed-separated, frozen-policy protocol defines the remaining
experiment required for a defensible performance claim.

\begin{thebibliography}{00}
\bibitem{mnih2015}
V. Mnih et al., ``Human-level control through deep reinforcement learning,''
\emph{Nature}, vol. 518, pp. 529--533, 2015.

\bibitem{hasselt2016}
H. van Hasselt, A. Guez, and D. Silver, ``Deep reinforcement learning with
Double Q-learning,'' in \emph{Proc. AAAI}, vol. 30, no. 1, 2016.

\bibitem{sutton2018}
R. S. Sutton and A. G. Barto, \emph{Reinforcement Learning: An Introduction},
2nd ed. Cambridge, MA, USA: MIT Press, 2018.

\bibitem{towers2024}
M. Towers et al., ``Gymnasium: A standard interface for reinforcement learning
environments,'' arXiv:2407.17032, 2024.

\bibitem{liu2022}
F. Liu et al., ``Integrated sensing and communications: Toward dual-functional
wireless networks for 6G and beyond,'' \emph{IEEE J. Sel. Areas Commun.},
vol. 40, no. 6, pp. 1728--1767, 2022.

\bibitem{riley2010}
G. F. Riley and T. R. Henderson, ``The ns-3 network simulator,'' in
\emph{Modeling and Tools for Network Simulation}. Berlin, Germany: Springer,
2010, pp. 15--34.
\end{thebibliography}

\end{document}
